% GENERATED FILE. Edit canonical lessons, then run npm run generate:books.
% canonical-source-hash: fnv1a64:671ed183bea1fd1e
% canonical-lessons: LA-C505-clemens, LA-C505-clementia, LA-C505-benignitas, LA-C505-benevolus, LA-C505-benevolentia

\chapter{Mercy and kindness}
\label{ch:la-mercy-and-kindness}
\hlchaptermodality{505}

\begin{chapteropening}
\textbf{By the end of this chapter:} \emph{I can talk about mercy, kindness and goodwill.}

Complete the last lesson of chapter 505: I can talk about mercy, kindness and goodwill.
\end{chapteropening}

\section[clēmēns, clēmentis]{clēmēns, clēmentis — merciful, mild}

\label{lesson:LA-C505-clemens}



\begin{quote}
Before the new one: say the Latin for to pretend, then the Latin for to hide.
\end{quote}

\subsection*{You'll want to know: clēmēns, clēmentis}
\textbf{clēmēns, clēmentis} — \textquotedblleft{}merciful, mild\textquotedblright{}.

\begin{quote}
Rēx clēmēns captīvīs ignōscit.
\end{quote}

\begin{quote}
\emph{The merciful king pardons the prisoners.}
\end{quote}

One ending; ablative usually \emph{clēmentī}. Of a ruler or a judge: gentle in punishing. English \emph{clement}.

\begin{cousinweb}
\textbf{clēmēns} is a native Latin adjective; Lewis \& Short calls its origin doubtful, so it is not settled. English \emph{clement} and \emph{inclement} come from it.
\end{cousinweb}

\begin{grammarlens}[title={What you've built}]
One more word for telling, explaining and describing at B1.
\end{grammarlens}

\subsection*{Your turn}
\begin{itemize}
\raggedright
  \item \emph{Say it:} \emph{clēmēns, clēmentis}
  \item \emph{Say it:} \emph{clēmēns, clēmentis}, once more
  \item \emph{Say it:} a sentence of your own with \emph{clēmēns, clēmentis}
  \item \emph{Recall:} say the Latin for to pretend, then the Latin for to hide, then say \emph{clēmēns, clēmentis} again
\end{itemize}

\begin{tcolorbox}[breakable,colback=teal!4,colframe=teal!35!black,title={Before you move on}]
What does clēmēns, clēmentis mean? (\textquotedblleft{}Merciful, mild\textquotedblright{}.) Say it once more.
\end{tcolorbox}
\section[clēmentia, clēmentiae]{clēmentia, clēmentiae — mercy, clemency}

\label{lesson:LA-C505-clementia}



\begin{quote}
Before the new one: say the Latin for to hide, then the Latin for merciful.
\end{quote}

\subsection*{You'll want to know: clēmentia, clēmentiae}
\textbf{clēmentia, clēmentiae} — \textquotedblleft{}mercy, clemency\textquotedblright{}.

\begin{quote}
Victor clēmentiam hostibus ostendit.
\end{quote}

\begin{quote}
\emph{The victor shows mercy to the enemy.}
\end{quote}

Feminine, first declension. A virtue of rulers: mercy toward those in your power.

\begin{cousinweb}
\textbf{clēmentia} is built on \emph{clēmēns} \textquotedblleft{}merciful\textquotedblright{} (Lewis \& Short). English \emph{clemency} comes from it.
\end{cousinweb}

\begin{grammarlens}[title={What you've built}]
One more word for telling, explaining and describing at B1.
\end{grammarlens}

\subsection*{Your turn}
\begin{itemize}
\raggedright
  \item \emph{Say it:} \emph{clēmentia, clēmentiae}
  \item \emph{Say it:} \emph{clēmentia, clēmentiae}, once more
  \item \emph{Say it:} a sentence of your own with \emph{clēmentia, clēmentiae}
  \item \emph{Recall:} say the Latin for to hide, then the Latin for merciful, then say \emph{clēmentia, clēmentiae} again
\end{itemize}

\begin{tcolorbox}[breakable,colback=teal!4,colframe=teal!35!black,title={Before you move on}]
What does clēmentia, clēmentiae mean? (\textquotedblleft{}Mercy, clemency\textquotedblright{}.) Say it once more.
\end{tcolorbox}
\section[benignitās, benignitātis]{benignitās, benignitātis — kindness, generosity}

\label{lesson:LA-C505-benignitas}



\begin{quote}
Before the new one: say the Latin for merciful, then the Latin for mercy.
\end{quote}

\subsection*{You'll want to know: benignitās, benignitātis}
\textbf{benignitās, benignitātis} — \textquotedblleft{}kindness, generosity\textquotedblright{}.

\begin{quote}
Benignitās hospitis nōs dēlectat.
\end{quote}

\begin{quote}
\emph{Our host's kindness delights us.}
\end{quote}

Feminine, third declension, from \emph{benignus} — kind. English \emph{benignity}.

\begin{cousinweb}
\textbf{benignitās} is built on \emph{benignus} \textquotedblleft{}kind\textquotedblright{}, which Lewis \& Short and EtymDB explain as \emph{bonus} or \emph{bene} \textquotedblleft{}good, well\textquotedblright{} with \emph{genus} \textquotedblleft{}kind\textquotedblright{}: of a good kind. English \emph{benign} comes from \emph{benignus}.
\end{cousinweb}

\begin{grammarlens}[title={What you've built}]
One more word for telling, explaining and describing at B1.
\end{grammarlens}

\subsection*{Your turn}
\begin{itemize}
\raggedright
  \item \emph{Say it:} \emph{benignitās, benignitātis}
  \item \emph{Say it:} \emph{benignitās, benignitātis}, once more
  \item \emph{Say it:} a sentence of your own with \emph{benignitās, benignitātis}
  \item \emph{Recall:} say the Latin for merciful, then the Latin for mercy, then say \emph{benignitās, benignitātis} again
\end{itemize}

\begin{tcolorbox}[breakable,colback=teal!4,colframe=teal!35!black,title={Before you move on}]
What does benignitās, benignitātis mean? (\textquotedblleft{}Kindness, generosity\textquotedblright{}.) Say it once more.
\end{tcolorbox}
\section[benevolus, benevola, benevolum]{benevolus, benevola, benevolum — kindly, well-wishing}

\label{lesson:LA-C505-benevolus}



\begin{quote}
Before the new one: say the Latin for mercy, then the Latin for kindness.
\end{quote}

\subsection*{You'll want to know: benevolus, benevola, benevolum}
\textbf{benevolus, benevola, benevolum} — \textquotedblleft{}kindly, well-wishing\textquotedblright{}.

\begin{quote}
Vīcīnus benevolus nōbīs semper adest.
\end{quote}

\begin{quote}
\emph{Our kindly neighbour is always there for us.}
\end{quote}

\emph{bene} 'well' + \emph{volō} 'to wish'. Often of a friend's goodwill.

\begin{cousinweb}
\textbf{benevolus} is built on \emph{bene} \textquotedblleft{}well\textquotedblright{} and \emph{volō} \textquotedblleft{}to wish\textquotedblright{} (Lewis \& Short).
\end{cousinweb}

\begin{grammarlens}[title={What you've built}]
One more word for telling, explaining and describing at B1.
\end{grammarlens}

\subsection*{Your turn}
\begin{itemize}
\raggedright
  \item \emph{Say it:} \emph{benevolus, benevola, benevolum}
  \item \emph{Say it:} \emph{benevolus, benevola, benevolum}, once more
  \item \emph{Say it:} a sentence of your own with \emph{benevolus, benevola, benevolum}
  \item \emph{Recall:} say the Latin for mercy, then the Latin for kindness, then say \emph{benevolus, benevola, benevolum} again
\end{itemize}

\begin{tcolorbox}[breakable,colback=teal!4,colframe=teal!35!black,title={Before you move on}]
What does benevolus, benevola, benevolum mean? (\textquotedblleft{}Kindly, well-wishing\textquotedblright{}.) Say it once more.
\end{tcolorbox}
\section[benevolentia, benevolentiae]{benevolentia, benevolentiae — goodwill, kindness}

\label{lesson:LA-C505-benevolentia}



\begin{quote}
Before the new one: say the Latin for kindness, then the Latin for kindly.
\end{quote}

\subsection*{You'll want to know: benevolentia, benevolentiae}
\textbf{benevolentia, benevolentiae} — \textquotedblleft{}goodwill, kindness\textquotedblright{}.

\begin{quote}
Sine benevolentiā nūlla amīcitia est.
\end{quote}

\begin{quote}
\emph{Without goodwill there is no friendship.}
\end{quote}

Feminine, first declension, from \emph{benevolus}. English \emph{benevolence}.

\begin{cousinweb}
\textbf{benevolentia} is built on \emph{benevolus} \textquotedblleft{}kindly, well-wishing\textquotedblright{} (Lewis \& Short). English \emph{benevolence} comes from it.
\end{cousinweb}

\begin{grammarlens}[title={What you've built}]
One more word for telling, explaining and describing at B1.
\end{grammarlens}

\subsection*{Your turn}
\begin{itemize}
\raggedright
  \item \emph{Say it:} \emph{benevolentia, benevolentiae}
  \item \emph{Say it:} \emph{benevolentia, benevolentiae}, once more
  \item \emph{Say it:} a sentence of your own with \emph{benevolentia, benevolentiae}
  \item \emph{Recall:} say the Latin for kindness, then the Latin for kindly, then say \emph{benevolentia, benevolentiae} again
\end{itemize}

\begin{tcolorbox}[breakable,colback=teal!4,colframe=teal!35!black,title={Before you move on}]
What does benevolentia, benevolentiae mean? (\textquotedblleft{}Goodwill, kindness\textquotedblright{}.) Say it once more.
\end{tcolorbox}
